% Capítulo incluido desde KalaScript.tex.
\section{Fundamentos de KalaScript}

\subsection{Qué quiero construir}

Con KalaScript quiero escribir ideas y conocimiento con una identidad visual propia.
Me interesa poder expresar matemáticas, computación, ciencias naturales,
filosofía, ciencias sociales, humanidades y otros campos. Construyo esa cobertura
gradualmente, a partir de conceptos y reglas que pueda reutilizar.

La idea nació de mi curiosidad por una escritura de ficción. A partir de esa
motivación, busco dar al proyecto una estructura concreta y documentada.

\textbf{Por ahora me concentro en la escritura:} no incorporo sonidos,
sílabas, pronunciaciones ni un alfabeto fonético. Mi objetivo es expresar ideas
directamente en KalaScript. Uso el español como apoyo explicativo en este manual;
las leyendas y los catálogos conservan sus funciones actuales.

\subsection{Cuatro elementos que distingo}

\begingroup\small
\begin{longtable}{@{}>{\raggedright\arraybackslash}p{\dimexpr 0.27\linewidth-1.0\tabcolsep\relax}>{\raggedright\arraybackslash}p{\dimexpr 0.73\linewidth-1.0\tabcolsep\relax}@{}}
\toprule
\textbf{Elemento} & \textbf{Función} \\ \midrule
\endfirsthead
\toprule
\textbf{Elemento} & \textbf{Función} \\ \midrule
\endhead
\bottomrule\endfoot
Significado & Qué concepto, objeto, relación o afirmación se expresa. \\[4pt]
Estructura & Qué modifica a qué, con qué argumentos y dentro de qué alcance. \\[4pt]
Escritura & Cómo se dibujan y distribuyen los signos. \\[4pt]
Registro digital & Cómo se conservan identificadores y expresiones para editarlos. \\[4pt]
\end{longtable}
\endgroup

Cambiar el dibujo de un signo no debe cambiar silenciosamente su significado. Cambiar un significado requiere una revisión explícita.

\subsection{Decisiones de escritura}

\begin{itemize}
\item Lectura principal de izquierda a derecha y bloques de arriba abajo.
\item Escritura monocroma: el color del atlas ayuda a orientarse, pero no significa nada.
\item Veintidós signos con categorías, operadores, marcadores y constructores definidos.
\item Índices decimales convencionales para referencias y relaciones. No se inventa todavía un sistema numeral.
\item Nombres convencionales de variables para la vinculación; fórmulas dentro de regiones matemáticas delimitadas.
\item Relaciones con argumentos ordenados y documentados; no se presupone que todo sea binario.
\item Alcance mediante grupos visibles; negación mediante un signo delante de un grupo.
\item Un pequeño fragmento de lógica clásica para los ejemplos de conjunción, negación, condicional y cuantificación.
\item Datos, ejemplos y figuras legibles sin servicios externos.
\end{itemize}

Los identificadores \texttt{K01}–\texttt{K22} son códigos de catálogo. No pertenecen al texto manuscrito ni tienen pronunciación.

\subsection{Distinciones obligatorias}

\textbf{Negar no es formar un antónimo.} «No caliente» puede incluir estados que no llamaríamos fríos. No comprar tampoco equivale a vender.

\textbf{Relación temporal no es causalidad.} Que un evento preceda a otro no afirma que lo cause.

\textbf{Parte y miembro son relaciones diferentes.} Una pieza de un objeto y un elemento de un conjunto no se expresan automáticamente de la misma manera.

\textbf{Fuente y certeza son diferentes.} Atribuir una afirmación a un informe no establece que sea verdadera ni asigna una probabilidad.

\textbf{Incertidumbre de medida no es error conocido.} Las cantidades y sus unidades deben acompañarse de la información necesaria para interpretar la medida.

\subsection{Qué sigue abierto}

Mantengo abiertos el diseño de las formas, el inventario de conceptos, posibles
modificadores gráficos, variantes manuscritas, abreviaturas y tipografía.
Tomo como referentes mecanismos de Hangul, Shavian, Blissymbolics, Gregg,
Dutton, Peirce, RDF e Iverson; evalúo cuáles encajan en el proyecto sin intentar
incorporarlos todos.
